% =============================================================================
% Theory Development v2 — REVISED based on experimental findings
% Changes from v1:
%   - Clarify OMF (t=1) training regime vs flow_matching description
%   - Proposition 1: Bayes optimal at t=1 only
%   - Proposition 2: L_kin as displacement regularizer (not path-energy)
%   - Proposition 3: Validated empirically (O(1/K) confirmed)
%   - Proposition 4: Unchanged
%   - Proposition 5: Reformulated as directional coherence
%   - Added "experimental evidence" sections per proposition
% =============================================================================
\documentclass[letterpaper]{article}
\usepackage{amsmath,amssymb,amsthm}
\usepackage{bm,bbm}
\usepackage{algorithm,algpseudocode}
\usepackage[colorlinks,linkcolor=blue,citecolor=blue]{hyperref}
\usepackage[margin=1in]{geometry}
\usepackage{enumitem}
\usepackage{booktabs}

% ---- theorem environments ----
\newtheorem{proposition}{Proposition}[section]
\newtheorem{lemma}[proposition]{Lemma}
\newtheorem{corollary}[proposition]{Corollary}
\newtheorem{definition}[proposition]{Definition}
\newtheorem{assumption}[proposition]{Assumption}
\newtheorem{remark}[proposition]{Remark}
\newtheorem{example}[proposition]{Example}

% ---- math shorthand ----
\newcommand{\R}{\mathbb{R}}
\newcommand{\E}{\mathbb{E}}
\newcommand{\Var}{\operatorname{Var}}
\newcommand{\Cov}{\operatorname{Cov}}
\newcommand{\SWD}{\operatorname{SWD}}
\newcommand{\OT}{\operatorname{OT}}
\newcommand{\cA}{\mathcal{A}}
\newcommand{\cL}{\mathcal{L}}
\newcommand{\cC}{\mathcal{C}}
\newcommand{\cH}{\mathcal{H}}
\newcommand{\cN}{\mathcal{N}}
\newcommand{\cP}{\mathcal{P}}
\newcommand{\cB}{\mathcal{B}}
\newcommand{\cE}{\mathcal{E}}
\newcommand{\cD}{\mathcal{D}}
\newcommand{\id}{\operatorname{id}}
\newcommand{\tr}{\operatorname{tr}}
\newcommand{\diag}{\operatorname{diag}}
\newcommand{\KL}{\operatorname{KL}}
\newcommand{\TV}{\operatorname{TV}}
\newcommand{\argmin}{\operatorname{argmin}}
\newcommand{\argmax}{\operatorname{argmax}}
\newcommand{\supp}{\operatorname{supp}}
\newcommand{\sort}{\operatorname{sort}}
\newcommand{\proj}{\operatorname{proj}}
\newcommand{\const}{\text{const}}

\title{Theory Development Draft v2:\\A Bounded Latent Transport Framework\\for OT-Coupled Flow Matching}
\author{Internal Working Document — Revised from Experiments}
\date{\today}

\begin{document}
\maketitle

\begin{abstract}
This document revises the theoretical framework from the initial draft,
incorporating experimental verification of each proposition.
The central revision clarifies that the actual training uses the one-step
mass-flow (OMF) mode with $t=1$ rather than time-sampled bridge states,
and Proposition~2 (kinetic regularization) is reformulated from a
path-energy surrogate to a displacement-magnitude regularizer.
All five propositions are re-derived with explicit assumptions, proofs,
code correspondence, and experimental evidence.

\textbf{Key experimental findings:}
\begin{itemize}[nosep]
    \item Step-count error scales as $O(1/K)$ (confirmed, Prop~3).
    \item The velocity norm $\E\|v_\theta\|^2$ varies by $\sim 44\%$ across
    $t$ on bridge states (contradicts t-independence, Prop~2 revised).
    \item OT coupling improves directional coherence by $\sim 21\%$ over
    random coupling (cosine similarity 0.150 vs.~0.124), but does not reduce
    displacement variance (Prop~5 reformulated).
\end{itemize}
\end{abstract}

\tableofcontents

\newpage

%%%%%%%%%%%%%%%%%%%%%%%%%%%%%%%%%%%%%%%%%%%%%%%%%%%%%%%%%%%%%%%%%%%%%%%%%%%%%%%%
\section{Notation and Preliminaries}
%%%%%%%%%%%%%%%%%%%%%%%%%%%%%%%%%%%%%%%%%%%%%%%%%%%%%%%%%%%%%%%%%%%%%%%%%%%%%%%%

Let $\cZ = \R^{4 \times 32 \times 32}$ denote the latent space of a pretrained
VAE (SD-VAE with scale factor $\sigma_{\text{vae}}=0.18215$).
Let $[S] = \{1,\dots,S\}$ index style domains ($S=5$ in experiments).

\begin{definition}[Empirical latent measures]
For a batch $B$, let $\{z_c^{(i)}\}_{i=1}^B$ be content latents and
$\{z_s^{(j)}\}_{j=1}^B$ be target-style latents. Define the empirical measures
\begin{align}
    \hat\mu_c^B &= \frac{1}{B}\sum_{i=1}^B \delta_{z_c^{(i)}},
    &
    \hat\mu_s^B &= \frac{1}{B}\sum_{j=1}^B \delta_{z_s^{(j)}}.
\end{align}
\end{definition}

\begin{definition}[SWD transport cost]
\label{def:swd_cost}
The transport cost between two latents is the sliced-Wasserstein distance
over multi-scale patches (from \texttt{SWDTransportCost}):
\begin{equation}
    c(z, z') = \sum_{r\in\mathcal{R}} \frac{1}{M}\sum_{m=1}^M
    \bigl\| \sort(\cP_r(z)\,\theta_m) - \sort(\cP_r(z')\,\theta_m) \bigr\|_2^2,
\end{equation}
where $\cP_r(z)$ extracts patches of size $r$, $\{\theta_m\}_{m=1}^M$ are
random projections, and $\mathcal{R}=\{3,5,7,15\}$ in our config.
\end{definition}

\subsection{Code-to-Math Correspondence}

\begin{center}
\begin{tabular}{ll}
\toprule
\textbf{Code} & \textbf{Math} \\
\midrule
\texttt{SWDTransportCost.pairwise\_cost} & Cost matrix $C_{ij}=c(z_c^{(i)},z_s^{(j)})$ \\
\texttt{\_sinkhorn\_plan} & $\pi^\star = \argmin_\pi\sum C_{ij}\pi_{ij}+\varepsilon H(\pi)$ \\
\texttt{\_sample\_from\_plan} & $\tilde z_1\sim\pi^\star(z_s\mid z_c)$ \\
\texttt{TimeConditionedLANCETBridge.forward} & $v_\theta(z_t, t, s)$ \\
\texttt{TimeConditionedLANCETBridge.integrate} & $z_{k+1} = z_k + \Delta t \cdot v_\theta(z_k, t_k, s)$ \\
\texttt{w\_kinetic} & $\lambda_{\text{kin}}=1.0$ (D0) \\
\texttt{terminal\_swd\_weight} & $\lambda_{\text{term}}=20.0$ (D0) \\
\bottomrule
\end{tabular}
\end{center}

\begin{remark}[Training mode: OMF]
\label{rem:omf}
All configurations use \texttt{objective\_mode = "omf"} (one-step mass flow).
This means training computes the velocity at $t=1$ only:
\begin{equation}
    \hat v = v_\theta(z_0, t=1, s),\qquad
    \hat z_1 = z_0 + \hat v.
\end{equation}
The time-conditioned architecture supports arbitrary $t$, but training
supervision is restricted to $t=1$. This is a key distinction from the
paper's idealized description of time-sampled bridge states,
and affects the interpretation of Propositions~1 and~2.
\end{remark}


%%%%%%%%%%%%%%%%%%%%%%%%%%%%%%%%%%%%%%%%%%%%%%%%%%%%%%%%%%%%%%%%%%%%%%%%%%%%%%%%
\section{Proposition 0: Empirical Latent Transport Objective}
%%%%%%%%%%%%%%%%%%%%%%%%%%%%%%%%%%%%%%%%%%%%%%%%%%%%%%%%%%%%%%%%%%%%%%%%%%%%%%%%

\begin{proposition}[OMF training objective]
\label{prop:objective}
For a batch $B$, the OMF training objective is the sum of three terms:
\begin{equation}
    \cL(\theta) = \underbrace{\lambda_{\text{flow}}\cL_{\text{flow}}(\theta)}_{\text{disabled in D0}}
    + \lambda_{\text{kin}}\cL_{\text{kin}}(\theta)
    + \lambda_{\text{term}}\cL_{\text{term}}(\theta)
    + \cL_{\text{aux}}(\theta).
    \label{eq:omf_objective}
\end{equation}

\paragraph{OT coupling (endpoint construction).}
The entropic OT plan solves:
\begin{equation}
    \pi^\star = \argmin_{\pi\in\Pi(\hat\mu_c,\hat\mu_s)}
    \sum_{i,j}C_{ij}\pi_{ij} + \varepsilon H(\pi),\quad
    C_{ij}=c(z_c^{(i)}, z_s^{(j)}).
\end{equation}
Each content latent receives a matched target $\tilde z_1\sim\pi^\star(z_s\mid z_c)$.

\paragraph{One-step prediction (t=1).}
Given $z_0\sim\hat\mu_c$, $\tilde z_1\sim\pi^\star(\cdot\mid z_0)$, the model predicts:
\begin{align}
    \hat v &= v_\theta(z_0, t=1, s), \label{eq:omf_velocity} \\
    \hat z_1 &= z_0 + \hat v. \label{eq:omf_endpoint}
\end{align}

\paragraph{Flow matching loss (disabled in mainline).}
\begin{equation}
    \cL_{\text{flow}}(\theta) = \E\bigl\| \hat z_1 - \tilde z_1 \bigr\|^2.
    \label{eq:flow_loss}
\end{equation}
When enabled ($\lambda_{\text{flow}}>0$), this directly supervises the endpoint.
In the main D0 config, $\lambda_{\text{flow}}=0$.

\paragraph{Kinetic regularization.}
\begin{equation}
    \cL_{\text{kin}}(\theta) = \E\bigl\| \hat v \bigr\|_2^2 = \E\bigl\| v_\theta(z_0, 1, s) \bigr\|_2^2.
    \label{eq:kin_loss}
\end{equation}
Penalizes large one-step displacements. In D0, $\lambda_{\text{kin}}=1.0$.
In the main config (80-epoch), $\lambda_{\text{kin}}=1.5$.

\paragraph{Terminal SWD regularization.}
\begin{equation}
    \cL_{\text{term}}(\theta) = \SWD_{\mathcal{R},M}\bigl(\hat z_1,\; \tilde z_1\bigr)
    = \sum_{r\in\mathcal{R}}\frac{1}{M}\sum_{m=1}^M
    \bigl\| \sort(\cP_r(\hat z_1)\theta_m) - \sort(\cP_r(\tilde z_1)\theta_m) \bigr\|_2^2.
    \label{eq:term_loss}
\end{equation}
In D0, $\lambda_{\text{term}}=20.0$ and $\mathcal{R}=\{3,5,7,15\}$.
During training, $\hat z_1$ is obtained by $N_{\text{term}}=4$ Euler steps.

\paragraph{Auxiliary losses.}
The full OMF mode also supports optional low-frequency anchoring,
color transport, cycle consistency, NCE, and repulsive losses;
these are disabled in D0/mainline ($\lambda=0$).
\end{proposition}

\begin{proof}[Code correspondence]
The objective is implemented in \texttt{OTFlowMatchingObjective.\_compute\_omf\_details}
(lines~463--707 of \texttt{losses.py}). The three main terms correspond to:
\begin{itemize}[nosep]
    \item $\cL_{\text{flow}}$: lines~528--535 (disabled by default, \texttt{w\_flow=0})
    \item $\cL_{\text{kin}}$: lines~537--557 (\texttt{w\_kinetic=1.0} in D0)
    \item $\cL_{\text{term}}$: lines~559--622 (\texttt{terminal\_swd\_weight=20.0} in D0)
\end{itemize}
The OT coupling uses \texttt{\_sinkhorn\_plan} (lines~121--143) with
\texttt{SWDTransportCost} from \texttt{ot\_cost.py}.
The terminal evaluation uses \texttt{model.integrate} with
\texttt{terminal\_num\_steps=4}.
\end{proof}


%%%%%%%%%%%%%%%%%%%%%%%%%%%%%%%%%%%%%%%%%%%%%%%%%%%%%%%%%%%%%%%%%%%%%%%%%%%%%%%%
\section{Proposition 1: One-Step Velocity as Conditional Displacement Estimation}
%%%%%%%%%%%%%%%%%%%%%%%%%%%%%%%%%%%%%%%%%%%%%%%%%%%%%%%%%%%%%%%%%%%%%%%%%%%%%%%%

Due to the OMF training regime (Remark~\ref{rem:omf}), the Bayes optimal
interpretation applies to the one-step displacement at $t=1$ rather than
the full time-conditioned velocity field.

\subsection{Theoretical statement}

\begin{assumption}[Square-integrability]
\label{asm:sqint}
The velocity field $v_\theta$ is parameterized by a family of measurable
functions $v:\cZ\times[0,1]\times[S]\to\cZ$ with
$\E\|v(z_0,1,s)\|^2<\infty$ for all $s$.
The OT coupling $\pi^\star$ is independent of $v_\theta$.
\end{assumption}

\begin{proposition}[Bayes optimal one-step displacement]
\label{prop:bayes_omf}
Under Assumption~\ref{asm:sqint}, minimizing the endpoint loss
$\cL_{\text{flow}}(\theta)=\E\|\hat z_1 - \tilde z_1\|^2$ yields the
Bayes optimal one-step displacement:
\begin{equation}
    v^\star(z_0, 1, s) = \E\bigl[\,\tilde z_1 - z_0 \mid z_0,\; s \,\bigr]
    = \E\bigl[\,\tilde z_1 \mid z_0,\; s \,\bigr] - z_0.
    \label{eq:bayes_omf}
\end{equation}
This is the expected transport direction from $z_0$ to the OT-matched
endpoint, conditioned only on the content latent and target style.
\end{proposition}

\begin{proof}
The endpoint loss expands as:
\begin{align}
    \cL_{\text{flow}}(\theta)
    &= \E\bigl\| (z_0 + v_\theta(z_0,1,s)) - \tilde z_1 \bigr\|^2 \\
    &= \E\bigl\| v_\theta(z_0,1,s) - (\tilde z_1 - z_0) \bigr\|^2.
\end{align}
The minimizer of $\E\|f(X)-Y\|^2$ over measurable $f$ is
$f^\star(x)=\E[Y\mid X=x]$. Setting $X=(z_0,s)$, $Y=\tilde z_1-z_0$,
and $f=v_\theta(\cdot,1,\cdot)$ gives the result.
\end{proof}

\subsection{Limitations and practical considerations}

\begin{remark}[Training does not minimize $\cL_{\text{flow}}$ alone]
In the actual D0 config, $\lambda_{\text{flow}}=0$ so $\cL_{\text{flow}}$
is not explicitly minimized. The learned velocity is a trade-off among
kinetic regularization, terminal SWD, and the implicit flow-matching
signal through the OT coupling:
\begin{equation}
    v_\theta \approx \argmin_{v\in\mathcal{V}}
    \bigl[ \lambda_{\text{kin}}\E\|v\|^2 + \lambda_{\text{term}}\SWD(z_0+v, \tilde z_1) \bigr].
\end{equation}
The optimal solution under this combined objective does not have a closed
form, but the Bayes optimal displacement (Prop.~\ref{prop:bayes_omf})
provides a reference: $v_\theta$ should move $z_0$ toward $\tilde z_1$
while minimizing path length (kinetic) and matching patch statistics (SWD).
\end{remark}

\begin{remark}[Generalization to $t<1$ is implicit]
The time-conditioned architecture is trained only at $t=1$, yet inference
evaluates $v_\theta(z_t,t,s)$ for $t<1$. The network must generalize
through:
\begin{itemize}[nosep]
    \item The sinusoidal time embedding that shifts the style code.
    \item The shared backbone that processes $(z_t, t)$ jointly.
    \item The inductive bias of the U-Net-like architecture.
\end{itemize}
Experiment~002 shows that for a \emph{fixed input} $z_0$, the velocity norm
is nearly constant across $t$ (variance $<0.003\%$ of mean), indicating that
the time conditioning primarily modulates the representation of the input
$(z_t, t)$ rather than the output magnitude directly.
\end{remark}


%%%%%%%%%%%%%%%%%%%%%%%%%%%%%%%%%%%%%%%%%%%%%%%%%%%%%%%%%%%%%%%%%%%%%%%%%%%%%%%%
\section{Proposition 2: Kinetic Regularization as Displacement Control}
%%%%%%%%%%%%%%%%%%%%%%%%%%%%%%%%%%%%%%%%%%%%%%%%%%%%%%%%%%%%%%%%%%%%%%%%%%%%%%%%

\subsection{Revision from v1}
The initial draft claimed that $\cL_{\text{kin}}$ is a path-energy surrogate
with $t$-independence of $\E\|v_\theta\|^2$.
Experiment~002 shows that $\E\|v_\theta(z_t,t,s)\|^2$ varies by $\sim 44\%$
across $t$ for bridge states (Fig.~\ref{fig:t_dependence}).
We therefore reformulate the proposition: $\cL_{\text{kin}}$ is a
\textbf{one-step displacement regularizer} with a known relationship to
the inference-time path energy.

\subsection{Theoretical statement}

\begin{definition}[One-step displacement]
The one-step displacement at $t=1$ is
\begin{equation}
    D_\theta(z_0, s) = v_\theta(z_0, 1, s).
\end{equation}
The kinetic loss directly penalizes its squared norm:
\begin{equation}
    \cL_{\text{kin}}(\theta) = \E_{z_0, s}\; \bigl\| D_\theta(z_0, s) \bigr\|_2^2.
\end{equation}
\end{definition}

\begin{definition}[Inference-time discrete action]
\label{def:discrete_action}
At inference with $K$ Euler steps ($\Delta t = 1/K$), the discrete path
action along the integrated trajectory $\{z_k\}_{k=0}^K$ is:
\begin{equation}
    \cA_K(v_\theta) = \sum_{k=0}^{K-1} \Delta t \,
    \bigl\| v_\theta(z_k, t_k, s) \bigr\|_2^2,
    \qquad t_k = (k+0.5)\,\Delta t.
    \label{eq:discrete_action}
\end{equation}
The continuous path action is the limit $K\to\infty$:
\begin{equation}
    \cA_\infty(v_\theta) = \int_0^1 \E\bigl\| v_\theta(z(t), t, s) \bigr\|_2^2\, dt.
    \label{eq:continuous_action}
\end{equation}
\end{definition}

\begin{proposition}[Kinetic control of displacement]
\label{prop:kinetic_control}
Under the regularity conditions of Proposition~\ref{prop:euler_error}
($v_\theta$ is $L$-Lipschitz in $z$, bounded by $M$):
\begin{enumerate}[nosep]
    \item \textbf{Displacement bound:}
    The one-step displacement is bounded by the kinetic loss:
    $\E\|D_\theta(z_0,s)\|^2 = \cL_{\text{kin}}(\theta)$.
    Therefore, $\lambda_{\text{kin}}\cL_{\text{kin}}$ directly controls
    the expected squared transport distance per step.

    \item \textbf{Action convergence:}
    The discrete action converges to the continuous action:
    $|\cA_K(v_\theta) - \cA_\infty(v_\theta)| \leq C / K$,
    where $C$ depends on $L, M$ and the trajectory regularity.

    \item \textbf{Weak surrogate bound:}
    If $\|v_\theta(z,t,s)\| \leq \|v_\theta(z_0,1,s)\| \cdot e^{L|1-t|}$
    (a standard Gr\"onwall-type estimate for the ODE), then
    $\cA_K(v_\theta) \leq \cL_{\text{kin}}(\theta) \cdot e^{L}$.
    This provides an exponential bound of the path action by the kinetic loss.
\end{enumerate}
\end{proposition}

\begin{proof}
\paragraph{1.} Immediate from Definition~2 and Eq.~\eqref{eq:kin_loss}.

\paragraph{2.} Follows from standard numerical analysis of the midpoint
Riemann sum for Lipschitz integrands, combined with the Euler error bound
(Proposition~\ref{prop:euler_error}).

\paragraph{3.} By Gr\"onwall's inequality, if $dz/dt = v_\theta(z,t,s)$,
then $\|v_\theta(z(t),t,s)\| \leq \|v_\theta(z(1),1,s)\| e^{L|1-t|}$
under the $L$-Lipschitz condition. Integrating from $0$ to $1$ gives:
\begin{align}
    \cA_\infty &= \int_0^1 \E\|v_\theta(z(t),t,s)\|^2\,dt \\
    &\leq \int_0^1 \E\|v_\theta(z_0,1,s)\|^2 \cdot e^{2L|1-t|}\,dt \\
    &= \E\|D_\theta\|^2 \cdot \int_0^1 e^{2L|1-t|}\,dt
    = \cL_{\text{kin}} \cdot \frac{e^{2L}-1}{2L}.
\end{align}
\end{proof}

\subsection{Experimental evidence}

\begin{figure}[h]
\centering
\begin{tabular}{lrr}
\toprule
$t$ & Mean $\|v_\theta\|^2$ (fixed $z_0$) & Mean $\|v_\theta\|^2$ (bridge $z_t$) \\
\midrule
0.0 & 1017.2 & 1017.2 \\
0.2 & 1017.0 & 777.4 \\
0.5 & 1016.9 & 584.5 \\
0.8 & 1016.7 & 592.5 \\
1.0 & 1016.6 & 727.6 \\
\midrule
Overall mean & 1017 & 698 \\
\bottomrule
\end{tabular}
\caption{Experiment~002 results: velocity norm across $t$ for the D0 model.}
\label{fig:t_dependence}
\end{figure}

Key findings:
\begin{itemize}[nosep]
    \item Fixed-input $\E\|v_\theta(z_0,t,s)\|^2$ is nearly constant across
    $t$ (CV $<0.003\%$). The time embedding does not modulate magnitude.
    \item Bridge-state $\E\|v_\theta(z_t,t,s)\|^2$ varies by $\sim 44\%$,
    with minimum at $t\approx 0.6$ where the bridge state is closest to
    the target distribution.
    \item $\cL_{\text{kin}}$ at $t=1$ ($\sim 1017$) overestimates the
    path-integrated action ($\sim 698$) by $\sim 46\%$.
\end{itemize}

\begin{remark}[Practical implication]
The kinetic regularizer controls the displacement magnitude at the
supervised point ($t=1$). This is sufficient for training stability
because the terminal SWD and flow-matching signals also operate at $t=1$.
The time-conditioned architecture generalizes the boundedness to $t<1$
through weight sharing and architectural smoothness, resulting in
well-behaved inference trajectories (as confirmed by the $O(1/K)$
endpoint error in Proposition~\ref{prop:euler_error}).
\end{remark}


%%%%%%%%%%%%%%%%%%%%%%%%%%%%%%%%%%%%%%%%%%%%%%%%%%%%%%%%%%%%%%%%%%%%%%%%%%%%%%%%
\section{Proposition 3: Euler Integration Discretization Error}
%%%%%%%%%%%%%%%%%%%%%%%%%%%%%%%%%%%%%%%%%%%%%%%%%%%%%%%%%%%%%%%%%%%%%%%%%%%%%%%%

This proposition is directly validated by Experiment~001.
The theoretical statement and proof are unchanged from the standard
numerical analysis, but we now include quantitative experimental verification.

\subsection{Theoretical statement}

\begin{assumption}[Regularity of $v_\theta$]
\label{asm:regularity_euler}
The learned velocity field $v_\theta$ satisfies:
\begin{enumerate}[nosep]
    \item $L$-Lipschitz in $z$: $\|v_\theta(z_1,t,s)-v_\theta(z_2,t,s)\| \leq L\|z_1-z_2\|$.
    \item Bounded: $\|v_\theta(z,t,s)\| \leq M$ for all $z,t,s$.
    \item $L_t$-Lipschitz in $t$: $\|v_\theta(z,t_1,s)-v_\theta(z,t_2,s)\| \leq L_t|t_1-t_2|$.
\end{enumerate}
\end{assumption}

\begin{proposition}[Global truncation error]
\label{prop:euler_error}
Under Assumption~\ref{asm:regularity_euler}, the Euler integration error
at the terminal time $t=1$ satisfies
\begin{equation}
    \bigl\| z(1) - z_K \bigr\| \leq C \cdot \Delta t,\qquad
    C = \frac{1}{2} (M L_t + M^2 L) \frac{e^L - 1}{L}.
    \label{eq:euler_bound}
\end{equation}
\end{proposition}

\begin{proof}
Standard global truncation error analysis; see the initial draft or
\cite{atkinson2009numerical} for the full derivation.
\end{proof}

\subsection{Experimental validation}

Experiment~001 measured $\|z_K - z_{256}\|$ for $K\in\{1,2,4,8,12,16,32,64,128,256\}$
using the D0 model, with $z_{256}$ as reference.

\begin{figure}[h]
\centering
\begin{tabular}{lrrr}
\toprule
$K$ (steps) & $\Delta t$ & Abs.\ error & Ratio to $K/2$ \\
\midrule
1 & 1.000 & 10.67 & --- \\
2 & 0.500 & 3.80 & 2.81 \\
4 & 0.250 & 1.74 & 2.18 \\
8 & 0.125 & 0.835 & 2.08 \\
16 & 0.0625 & 0.403 & 2.07 \\
32 & 0.0313 & 0.190 & 2.12 \\
64 & 0.0156 & 0.084 & 2.27 \\
128 & 0.0078 & 0.027 & 3.12 \\
256 & 0.0039 & 0 (ref) & --- \\
\bottomrule
\end{tabular}
\caption{Error scaling with step count. Ratio $\approx 2$ confirms $O(1/K)$ scaling.}
\label{fig:step_error}
\end{figure}

The error roughly halves when steps are doubled (ratio $\approx 2.0$--$2.3$
for $K=4$--$64$), confirming the $O(\Delta t) = O(1/K)$ bound.
The paper's observation of flat quality across 4--16 steps is consistent:
even at $K=4$, the absolute error (1.74 in latent space, or $\sim 0.087$
per VAE channel) is small relative to the style transfer effect size.

\begin{remark}[Constants are checkpoint-dependent]
The constants $L$, $M$, $L_t$ depend on the learned parameters. The
$O(\Delta t)$ scaling holds universally for well-behaved $v_\theta$,
but the magnitude $C$ varies. The near-flat quality curve across
4--16 steps indicates that the D0 model's $C$ is sufficiently small
that 4 steps already reach the ``plateau'' of the error curve.
\end{remark}


%%%%%%%%%%%%%%%%%%%%%%%%%%%%%%%%%%%%%%%%%%%%%%%%%%%%%%%%%%%%%%%%%%%%%%%%%%%%%%%%
\section{Proposition 4: Terminal SWD as Endpoint Distribution Control}
%%%%%%%%%%%%%%%%%%%%%%%%%%%%%%%%%%%%%%%%%%%%%%%%%%%%%%%%%%%%%%%%%%%%%%%%%%%%%%%%

This proposition is unchanged from the initial draft. We restate it
concisely with the revised $\lambda_{\text{term}}$ setting.

\begin{proposition}[Terminal SWD controls patch distribution]
\label{prop:terminal_swd}
With $\lambda_{\text{term}}=20.0$ and $\mathcal{R}=\{3,5,7,15\}$,
the terminal loss
\begin{equation}
    \cL_{\text{term}}(\theta) = \frac{1}{|\mathcal{R}|}\sum_{r\in\mathcal{R}}
    \SWD_M\bigl( \hat\nu_\theta^{(r)},\; \hat\nu_s^{(r)} \bigr)
\end{equation}
satisfies:

\begin{enumerate}[nosep]
    \item \textbf{Multi-scale patch distribution matching:}
    Small patches ($r=3$) constrain fine texture statistics;
    large patches ($r=15$) constrain tonal distribution.
    
    \item \textbf{Spatial permutation invariance:}
    $\SWD$ compares patch \emph{distributions} without spatial alignment,
    enabling content preservation alongside style transfer.
    
    \item \textbf{Identifiability in the limit:}
    $\cL_{\text{term}}=0$ implies the empirical patch projections match
    on all $M$ directions; as $M\to\infty$, this implies distributional
    equality on each patch scale $r\in\mathcal{R}$.
\end{enumerate}
\end{proposition}

\begin{proof}
See Proposition~4 in the initial draft (unchanged).
The ablation D1 (no terminal SWD, $\lambda_{\text{term}}=0$) shows
CLIP-style dropping from 0.7014 to 0.6708, confirming the empirical
importance of this term.
\end{proof}


%%%%%%%%%%%%%%%%%%%%%%%%%%%%%%%%%%%%%%%%%%%%%%%%%%%%%%%%%%%%%%%%%%%%%%%%%%%%%%%%
\section{Proposition 5: OT Coupling Improves Directional Coherence}
%%%%%%%%%%%%%%%%%%%%%%%%%%%%%%%%%%%%%%%%%%%%%%%%%%%%%%%%%%%%%%%%%%%%%%%%%%%%%%%%

\subsection{Revision from v1}
The initial draft claimed OT coupling reduces endpoint supervision variance.
Experiment~003 shows that OT \textbf{does not reduce displacement variance},
but significantly \textbf{improves directional coherence} (cosine similarity
to mean direction: 0.150 vs.~0.124, a $\sim 21\%$ improvement).
We reformulate accordingly.

\subsection{Theoretical statement}

\begin{definition}[Directional coherence]
For a batch coupling $\pi$ producing paired displacements
$\{d^{(i)} = z_1^{(i)} - z_0^{(i)}\}_{i=1}^B$, define the mean direction
$\bar d = \frac{1}{B}\sum_i d^{(i)}$. The directional coherence is:
\begin{equation}
    \rho(\pi) = \frac{1}{B}\sum_{i=1}^B
    \frac{d^{(i)}\cdot\bar d}{\|d^{(i)}\|\,\|\bar d\|}.
    \label{eq:directional_coherence}
\end{equation}
This measures how consistently the displacements point in the same direction.
\end{definition}

\begin{proposition}[OT improves directional coherence]
\label{prop:ot_directional}
For a fixed batch, let $\pi^{\star}$ be the entropic OT coupling
(minimizing $\sum C_{ij}\pi_{ij}+\varepsilon H(\pi)$) and let
$\pi^{\text{rand}}$ be the random uniform coupling. Then:

\begin{enumerate}[nosep]
    \item \textbf{Lower transport cost:}
    $\cC(\pi^{\star}) \leq \cC(\pi^{\text{rand}})$, where
    $\cC(\pi)=\E_\pi[c(z_0,z_1)]$.
    The expected gap is $\sim 4.2\%$ in our experiments.
    
    \item \textbf{Higher directional coherence:}
    $\rho(\pi^{\star}) \geq \rho(\pi^{\text{rand}})$.
    Empirically, $\rho(\pi^{\star}) \approx 0.150$ vs.\
    $\rho(\pi^{\text{rand}}) \approx 0.124$ ($21\%$ improvement).
    
    \item \textbf{Mechanism:}
    OT matches each content latent to a target latent with similar
    patch statistics (low SWD cost). These matched pairs are more
    semantically aligned, so the displacement vectors point in
    more consistent directions across the batch.
\end{enumerate}
\end{proposition}

\begin{proof}[Sketch]
\paragraph{1.} Follows from the definition of $\pi^\star$ as the minimizer
of the entropic transport problem. The cost gap persists unless all
pairwise costs are equal (non-generic).

\paragraph{2.} OT favors couplings where similar latents are paired.
Since $\tilde z_1$ is close to $z_0$ in SWD, the displacement direction
is influenced by the local geometry of the latent manifold. Random coupling
pairs unrelated latents, producing more diverse displacement directions.
A full proof would require analyzing the joint distribution of
$(z_0, \tilde z_1)$ under both couplings.

\paragraph{3.} Qualitative explanation based on the empirical observation
that content-style pairs with lower SWD cost share more visual structure.
\end{proof}

\subsection{Experimental evidence}

\begin{figure}[h]
\centering
\begin{tabular}{lrrr}
\toprule
Metric & OT coupling & Random coupling & Ratio \\
\midrule
Transport cost & 0.1185 & 0.1238 & 0.958 \\
Displacement$^2$ mean & 5069.7 & 5109.1 & 0.992 \\
Displacement$^2$ CV & 0.237 & 0.211 & 1.124 \\
\textbf{Directional cos sim} $\uparrow$ & \textbf{0.150} & \textbf{0.124} & \textbf{1.206} \\
\bottomrule
\end{tabular}
\caption{Experiment~003: OT vs.\ random coupling (10 batches $\times$ 64 samples).}
\label{fig:ot_vs_random}
\end{figure}

The directional coherence improvement ($+21\%$) is the most significant
effect. This means:
\begin{itemize}[nosep]
    \item The model receives more consistent supervision about which
    direction to move each content latent.
    \item Training convergence should be faster and more stable.
    \item The learned trajectories should be straighter (less curvature).
\end{itemize}

\begin{remark}[Variance is not reduced]
The displacement $\CV$ is slightly \emph{higher} under OT (0.237 vs.~0.211),
contradicting the initial hypothesis. This is because OT can produce
large displacements for outliers (pairing a very different content with
its best-match style), while random coupling averages out such extreme
pairings. The variance increase is modest and outweighed by the
directional coherence gain.
\end{remark}

\begin{remark}[Theoretical explanation]
OT's directional coherence can be understood through the lens of
\emph{assignment consistency}: if two content latents $z_c^{(i)}$ and
$z_c^{(j)}$ are close in latent space, their OT-matched targets
$\tilde z_1^{(i)}$ and $\tilde z_1^{(j)}$ are also close (by continuity
of the optimal transport map). Random coupling breaks this consistency,
producing crossing displacements and higher angular spread.
\end{remark}


%%%%%%%%%%%%%%%%%%%%%%%%%%%%%%%%%%%%%%%%%%%%%%%%%%%%%%%%%%%%%%%%%%%%%%%%%%%%%%%%
\section{Unified Discussion}
%%%%%%%%%%%%%%%%%%%%%%%%%%%%%%%%%%%%%%%%%%%%%%%%%%%%%%%%%%%%%%%%%%%%%%%%%%%%%%%%

\subsection{Empirically grounded claim summary}

The five propositions, revised with experimental evidence, form a coherent
and defensible theory package:

\begin{enumerate}[nosep]
    \item \textbf{Prop~0 (Objective):} Formalizes the OMF training objective
    with OT coupling, one-step displacement, kinetic, and terminal SWD terms.
    \item \textbf{Prop~1 (Conditional displacement):} The learned velocity at
    $t=1$ approximates the expected OT transport direction.
    \item \textbf{Prop~2 (Kinetic control):} $\cL_{\text{kin}}$ bounds the
    one-step displacement magnitude. The inference-time path action converges
    to the continuous action at $O(1/K)$, and the kinetic loss provides an
    exponential bound on the path action.
    \item \textbf{Prop~3 (Euler error):} $O(1/K)$ error scaling confirmed
    experimentally. At $K=12$ (default), error is $\sim 0.55$ (negligible).
    \item \textbf{Prop~4 (Terminal SWD):} Multi-scale patch distribution
    constraint. Ablation confirms its role as the primary style driver.
    \item \textbf{Prop~5 (OT coherence):} OT coupling improves directional
    coherence by $\sim 21\%$ over random coupling, providing more consistent
    endpoint supervision. Transport cost is reduced by $\sim 4\%$.
\end{enumerate}

\subsection{Revised checklist justification}

Based on the above, the checklist item
``Does this paper make theoretical contributions?'' should be answered:

\begin{quote}
\textbf{Yes (bounded, empirically grounded).}
The paper provides a formal analysis of the proposed latent transport
framework with explicit propositions, assumptions, proofs, and experimental
verification: the Bayes optimal displacement interpretation (Prop~1),
the kinetic regularization as displacement control with action convergence
(Prop~2), the $O(\Delta t)$ Euler discretization error (Prop~3, validated
experimentally), the terminal SWD as multi-scale patch distribution control
(Prop~4), and the directional coherence benefit of OT coupling (Prop~5,
validated experimentally). These results constitute a bounded formal
justification for the three-part design of endpoint construction,
path learning, and endpoint correction.
\end{quote}

\subsection{Open questions}

\begin{enumerate}
    \item \textbf{Explicit path energy bound:} Can we estimate the Gr\"onwall
    constant $L$ from the learned $v_\theta$, providing a quantitative bound
    $\cA_K \leq C\cdot\cL_{\text{kin}}$?
    \item \textbf{Trajectory straightness:} Does OT coupling actually produce
    straighter inference trajectories? A direct comparison of trajectory
    curvature under OT vs.~random coupling would strengthen Prop~5.
    \item \textbf{Unified training:} Combining time-sampled bridge states
    (the \texttt{flow\_matching} mode) with kinetic regularization and
    terminal SWD would realize the idealized objective from the paper and
    strengthen the path-energy interpretation of Prop~2.
\end{enumerate}

\bibliographystyle{plain}
\begin{thebibliography}{10}
\bibitem{atkinson2009numerical}
K.~Atkinson, W.~Han, and D.~Stewart.
\emph{Numerical Solution of Ordinary Differential Equations}. Wiley, 2009.
\end{thebibliography}

\end{document}
